\expandafter\def\csname HistoricalNoteText16\endcsname{[Handwritten remark.] On 31 Dec. already almost 11 degrees. Compare pp. 581 and 647. [P. 203 of this volume.]}
\clearpage\SourceLeaf{204}%
\pdfbookmark[1]{Further Report. Ceres. December 1801}{ceres180112}
\fancyhead[L]{\small [CERES.]}
\begin{center}
{\small FURTHER REPORT ON THE LONG-SUSPECTED NEW MAJOR PLANET OF OUR SOLAR SYSTEM.\par}
[CERES.]\par
\medskip\hrule\medskip
{\small Monatliche Correspondenz zur Beförderung der Erd- und Himmels-Kunde, edited by Baron \textsc{von Zach}. Vol. III. December 1801.\par}
\medskip\hrule
\end{center}\medskip
\noindent —\quad —\quad —\par
The investigations and calculations that Dr.\ \textsc{Gauss} in Brunswick recently shared with us give us great hope of help and relief. At the same time they give us new grounds for thinking it highly probable that the new object discovered by \textsc{Piazzi} is indeed a planetary body moving between the orbits of Mars and Jupiter according to \textsc{Kepler}'s laws.

We are all the more eager to share his calculations because his new elliptical orbit differs considerably from the elliptical orbit of Dr.\ \textsc{Burckhardt} and the two circular orbits of Dr.\ \textsc{Olbers} and Father \textsc{Piazzi}, which we presented in previous issues along with the positions of the planet predicted from them. This month their positions may differ from \textsc{Gauss}'s by as much as 6 to 7 degrees in geocentric longitude. It is therefore important to pass these remarks on to practical astronomers as soon as possible. They will then learn that they \textit{must} extend the region of the sky in which they are to look for this new and elusive object by another 6 to 7 degrees to the east. For 1) according to Dr.\ \textsc{Gauss}'s calculation, \textsc{Piazzi}'s observations lie almost midway between perihelion and aphelion, whereas Drs.\ \textsc{Burckhardt} and \textsc{Olbers} had supposed them to lie near one or the other.

2) According to his calculation the eccentricity of the orbit is more appreciable than Father \textsc{Piazzi} believes. The longitudes predicted from his elliptical elements can therefore easily differ by 7 degrees\HistoricalNote{16} this month from those calculated under the circular hypotheses.

3) \textsc{Gauss}'s elements fully vindicate \textsc{Piazzi}'s observations and entirely clear them of the suspicion of inaccuracy that Drs.\ \textsc{Burckhardt} and \textsc{Olbers} had entertained when testing the observations by differences.

\textsc{Gauss}'s ellipse demonstrates all this. Astronomers will see from the accuracy with which it represents all of \textsc{Piazzi}'s observations how much confidence it must inspire. Dr.\ \textsc{Gauss} was led to these calculations by some investigations in physical astronomy
\SourceLeaf{205}%
that brought him to some additions of no small importance to the theory of determining the motion of celestial bodies in conic sections of every kind. He was kind enough to share some of these with us. We hope to entertain our astronomical readers with them on another occasion, since an account of them here would take us too far from our subject. We therefore confine ourselves to what bears directly on the investigation of the orbit of \textit{Ceres Ferdinandea}.

Dr.\ \textsc{Gauss} first chose the three observations of 2 and 22 January and 11 February to determine the orbit. He took these data exactly as given on p.\ 280 of the September issue of the M.\ C. Using a procedure of his own, he found the following elements as early as the fourth trial:
\begin{center}
{\small\setlength{\tabcolsep}{2pt}
\begin{tabular}{l r}
Aphelion $\ .\ .\ .\ .\ .\ .\ .\ .\ .\ .$&$330^0\ 14'\ 33''$\\
$\AscendingNode\ .\ .\ .\ .\ .\ .\ .\ .\ .\ .\ .\ .\ .$&$81\quad8\ 50$\\
Orbital inclination $\ .\ .\ .\ .\ .\ .\ .\ .$&$10\ 32\ 19$\\
Logarithm of the semimajor axis $\ .\ .$&$0.4381058$\\
Eccentricity $\ .\ .\ .\ .\ .\ .\ .\ .\ .\ .$&$0.0832836$\\
\multicolumn{2}{l}{Epoch: 31 Dec. 1800 at mean noon in Palermo}\\
mean heliocentric longitude: $\ .\ .\ .$&$77^0\ 54'\ 29''$\\[4pt]
\multicolumn{2}{c}{Hence:}\\[3pt]
greatest equation of the centre $\ .\ .\ .$&$9^0\ 32'\ 57''$\\
semimajor axis $\ .\ .\ .\ .\ .\ .\ .\ .$&$2.74226$\\
sidereal period $\ .\ .\ .\ .\ .\ .$&$1658\frac23$ days\\
mean daily sidereal motion&$781''.355$
\end{tabular}}
\end{center}
These elements represent the two outer observations exactly and the middle one with an error of $2''$ in longitude and latitude. Due account was taken of aberration and precession.

After this first successful trial Dr.\ \textsc{Gauss} undertook a second calculation of the orbit. He left out the observations of 2 and 22 January and chose those of 1 and 21 January instead, along with the observation of 11 February, which he wished to keep so that the interval would be as long as possible. On the fourth hypothesis this calculation gave him the following elements:
\begin{center}
{\small\setlength{\tabcolsep}{2pt}
\begin{tabular}{l r}
Aphelion $\ .\ .\ .\ .\ .\ .\ .\ .\ .$&$330^0\ 33'\ 20''$\\
$\AscendingNode\ .\ .\ .\ .\ .\ .\ .\ .\ .\ .\ .\ .\ .$&$81\quad2\ 35$\\
Orbital inclination $\ .\ .\ .\ .\ .\ .\ .$&$10\ 36\ 30$\\
Logarithm of the semimajor axis&$0.4370335$\\
Eccentricity $\ .\ .\ .\ .\ .\ .\ .\ .\ .$&$0.0705553$\\
Epoch as above $\ .\ .\ .\ .\ .\ .\ .$&$76^0\ 28'\ 14''.27$\\[4pt]
\multicolumn{2}{c}{Hence:}\\[3pt]
greatest equation of the centre $\ .\ .\ .$&$8^0\ 5'\ 19''$\\
semimajor axis $\ .\ .\ .\ .\ .\ .\ .\ .$&$2.73548$\\
sidereal period $\ .\ .\ .\ .\ .\ .$&$1652\frac12$ days\\
mean daily sidereal motion $\ .$&$784''.25$
\end{tabular}}
\end{center}
\Needspace{16\baselineskip}
With these elements all of \textsc{Piazzi}'s observations listed on p.\ 280 of the September issue agree as follows:
\begin{center}
{\small\setlength{\tabcolsep}{4pt}
\begin{tabular}{r|r|r|r|r}
1801&\multicolumn{2}{c|}{Calculated}&\multicolumn{2}{c}{Errors in}\\
&Longitude&Latitude&Longitude&Latitude\\\hline
January 1&$53^0\ 22'\ 58''42$&$3^0\ 6'\ 42''09$&$+0''12$&$-0''01$\\
2&$53\ 19\ 37.02$&$3\ 2\ 23.78$&$-7.28$&$-1.12$\\
$\mathord{:}\mathord{:}$ 3&$53\ 16\ 43.67$&$2\ 58\ 6.70$&$-14.93$&$-3.20$\\
4&$53\ 14\ 14.03$&$2\ 53\ 51.16$&$-1.47$&$-4.44$\\
$\mathord{:}\mathord{:}$ 10&$53\ 7\ 54.51$&$2\ 28\ 53.13$&$-4.59$&$-7.47$\\
13&$53\ 10\ 18.77$&$2\ 16\ 48.78$&$-18.83$&$-10.92$\\
14&$53\ 11\ 55.25$&$2\ 12\ 51.23$&$-5.95$&$-5.47$\\
19&$53\ 26\ 0.37$&$1\ 53\ 34.19$&$+1.17$&$-4.01$\\
21&$53\ 34\ 22.68$&$1\ 46\ 5.98$&$+1.38$&$-0.02$\\
22&$53\ 39\ 7.88$&$1\ 42\ 25.06$&$+6.08$&$-3.04$
\end{tabular}
\par\medskip
\begin{tabular}{r|r|r|r|r}
1801&\multicolumn{2}{c|}{Calculated}&\multicolumn{2}{c}{Errors in}\\
&Longitude&Latitude&Longitude&Latitude\\\hline
January 23&$53^0\ 44'\ 15''74$&$1^0\ 38'\ 46''25$&$+0''04$&$-5''85$\\
28&$54\ 15\ 20.88$&$1\ 21\ 4.07$&$+5.18$&$-2.83$\\
30&$54\ 30\ 4.25$&$1\ 14\ 13.73$&$+5.25$&$-2.27$\\
31&$54\ 38\ 11.25$&$1\ 10\ 51.99$&$+3.95$&$-2.61$\\
$\mathord{:}\mathord{:}$ February 1&$54\ 46\ 28.37$&$1\ 7\ 32.09$&$+9.07$&$+1.19$\\
2&$54\ 55\ 5.09$&$1\ 4\ 14.25$&$+7.19$&$+3.75$\\
5&$55\ 22\ 50.25$&$0\ 54\ 32.88$&$+6.85$&$+3.98$\\
8&$55\ 53\ 23.15$&$0\ 45\ 9.20$&$-6.38$&$+4.20$\\
11&$56\ 26\ 39.97$&$0\ 36\ 2.90$&$-0.03$&$+0.00$
\end{tabular}}
\end{center}
\SourceLeaf{206}%
Excellent as this agreement is, Dr.\ \textsc{Gauss} himself considers it quite possible that his elements may depart considerably more from the true ones than they differ among themselves, since the portion of the orbit traversed is so small and amounts to only $9^0\ 15'\ 35''$ from the first to the last observation. Yet since this orbit is \textit{possible}, if not \textit{highly probable}, practical astronomers have doubly good reason to take it into account in looking for this object and to choose it as their principal guide. None of the orbits tried and calculated so far brings all the observations into such close and exact agreement as this ellipse of \textsc{Gauss}.

We should also mention that in all these calculations Dr.\ \textsc{Gauss} allowed both for the one-minute misprint in the position of the Sun pointed out on p.\ 365 of the October issue and for a second misprint on 13 Jan., likewise in the longitude of the Sun. He took the latter longitude to be $9^{\mathrm z}\ 23^0\ 13'\ 13.''8$.

This was how far Dr.\ \textsc{Gauss} had got in calculating the orbit of \textit{Ceres Ferdinandea} when he was kind enough to share his calculations with us. In the meantime we had received \textsc{Piazzi}'s paper together with his corrected observations, and we gave a faithful extract in the November issue. These changes were slight, the most substantial being the reduction of $15''$ in the right ascension for 11 Feb., and they made little difference to the preceding elements. Yet some more significant printing and calculation errors had crept into the reduction of these observations here and there. Since these errors could readily be removed and the elements soon improved accordingly, we sent Dr.\ \textsc{Gauss} all of \textsc{Piazzi}'s observations in the original Italian publication for this purpose. But he had already started to calculate a new orbit before receiving our letter. Slight as the deviations of his second set of elements from the observations already were, they still displayed a striking regularity, from which it could be predicted with certainty that the errors could be reduced considerably further. He tried it, and when he received the corrected observations that had been sent to him, as reduced and calculated by \textsc{Piazzi} himself, he had already done most of the work. He thought it best to finish it first, especially since he knew in advance that the longitude for 11 Feb. given by the new elements would be $6''$ smaller and would thus of itself come almost halfway toward the corrected observed value, which of course was still unknown to him. In this way he found the following \textit{third} set of elements:
\begin{center}
{\small\setlength{\tabcolsep}{2pt}
\begin{tabular}{l r}
Aphelion $\ .\ .\ .\ .\ .\ .\ .\ .\ .\ .$&$326^0\ 53'\ 50''$\\
$\AscendingNode\ .\ .\ .\ .\ .\ .\ .\ .\ .\ .\ .\ .\ .$&$81\quad1\ 44$\\
Orbital inclination $\ .\ .\ .\ .\ .\ .\ .\ .$&$10\ 36\ 21$\\
Logarithm of the semimajor axis $\ .$&$0.4414902$\\
Eccentricity $\ .\ .\ .\ .\ .\ .\ .\ .\ .\ .$&$0.0819603$\\
Epoch $\ .\ .\ .\ .\ .\ .\ .\ .\ .\ .\ .\ .$&$77^0\ 34'\ 28''$\\[4pt]
\multicolumn{2}{c}{Hence also:}\\[3pt]
Greatest equation of the centre $\ .\ .\ .$&$9^0\ 23'\ 57''$\\
Mean daily sidereal heliocentric motion&$772''275$\\
Daily tropical heliocentric motion $\ .$&$772.413$\\
Tropical period $\ .\ .\ .\ .\ .\ .$&$1677.8$ days
\end{tabular}}
\end{center}
These elements represent \textsc{Piazzi}'s observations, as reduced by us and listed on p.\ 280 of the September issue, as follows:
\expandafter\def\csname HistoricalNoteText17\endcsname{In \textsc{Piazzi}'s original observations on p.\ 565 of the November issue, the longitudes of the Sun on 13 January and 8 February are one minute too large. All geocentric latitudes are about $10''$ too small. This is undoubtedly because \textsc{Piazzi} used a different obliquity of the ecliptic in reducing his observations, probably that of $23^0\ 28'\ 1''$ from the \textit{Conn.\ d.\ tems Année} IX.}
\SourceLeaf{207}%
\begin{center}
{\small\setlength{\tabcolsep}{4pt}
\begin{tabular}{r|r|r|r|r}
1801&\multicolumn{2}{c|}{Calculated}&\multicolumn{2}{c}{Errors in}\\
&Longitude&Latitude&Longitude&Latitude\\\hline
January 1&$53^0\ 23'\ 2''34$&$3^0\ 6'\ 43''63$&$+4''04$&$+1''53$\\
2&$53\ 19\ 41.24$&$3\ 2\ 25.68$&$-3.06$&$+0.78$\\
3&$53\ 16\ 48.05$&$2\ 58\ 8.97$&$-10.35$&$-0.93$\\
4&$53\ 14\ 18.47$&$2\ 53\ 53.79$&$+2.97$&$-1.81$\\
10&$53\ 7\ 58.37$&$2\ 28\ 57.12$&$-0.73$&$-3.48$\\
13&$53\ 10\ 21.60$&$2\ 16\ 52.89$&$-16.00$&$-6.81$\\
14&$53\ 11\ 57.70$&$2\ 12\ 55.36$&$-3.50$&$-1.34$\\
19&$53\ 26\ 0.59$&$1\ 53\ 38.01$&$+1.39$&$-0.19$\\
21&$53\ 34\ 21.99$&$1\ 46\ 9.53$&$+0.69$&$+3.53$\\
22&$53\ 39\ 6.69$&$1\ 42\ 28.45$&$+4.89$&$+0.35$
\end{tabular}
\par\medskip
\begin{tabular}{r|r|r|r|r}
1801&\multicolumn{2}{c|}{Calculated}&\multicolumn{2}{c}{Errors in}\\
&Longitude&Latitude&Longitude&Latitude\\\hline
January 23&$53^0\ 44'\ 14''08$&$1^0\ 38'\ 49''44$&$-1''62$&$-2''66$\\
28&$54\ 15\ 17.11$&$1\ 21\ 5.91$&$+1.41$&$-0.99$\\
30&$54\ 30\ 9.76$&$1\ 14\ 15.12$&$+0.76$&$-0.88$\\
31&$54\ 38\ 6.44$&$1\ 10\ 52.81$&$-0.86$&$-1.79$\\
February 1&$54\ 46\ 23.22$&$1\ 7\ 32.54$&$+3.92$&$+1.64$\\
2&$54\ 54\ 59.71$&$1\ 4\ 14.30$&$+1.81$&$+3.80$\\
5&$55\ 22\ 44.30$&$0\ 54\ 31.72$&$+0.90$&$+2.82$\\
8&$55\ 53\ 17.01$&$0\ 45\ 6.65$&$-12.49$&$+1.63$\\
11&$56\ 26\ 34.10$&$0\ 35\ 58.96$&$-5.90$&$-3.94$
\end{tabular}}
\end{center}
The errors and corrections indicated by \textsc{Piazzi} now give the following changes. With the right ascension of 11 Feb. reduced by $15''$ and with $23^0\ 28'\ 12''$ for the obliquity of the ecliptic, the longitude becomes $56^0\ 26'\ 26''1$ and the latitude $35'\ 59''7$. Thus the error in longitude is $+8''0$ and the error in latitude $-0''74$. In recalculating the reduction of right ascension and declination for those observations where \textsc{Piazzi}'s figures differ considerably from ours, Dr.\ \textsc{Gauss} found the longitude on 8 Feb. to be $55^0\ 53'\ 17''7$. The error is therefore $-0''7$\HistoricalNote{17}. For the rest, and even for the longitude of 13 January, which differs so greatly in \textsc{Piazzi}, Dr.\ \textsc{Gauss} obtains the same results as we do. This particular observation deviates the most from both these latest elements and the preceding ones, and it cannot be represented exactly without forcing the other observations. Dr.\ \textsc{Gauss} therefore suspects that some small oversight occurred in connection with it. It is easy to see, moreover, that the other small corrections given by \textsc{Piazzi} (November issue, p.\ 573) make little difference to the accuracy with which the observations are represented by these latest elements.

Slight as all the errors are, Dr.\ \textsc{Gauss} nevertheless made a fourth attempt to increase the accuracy a little further. With these latest elements all the longitude errors in February are positive if \textsc{Piazzi}'s small corrections are also taken into account. Toward the end the latitude errors also tend mostly to one side. Dr.\ \textsc{Gauss} therefore sought new elements for which the calculated longitudes and latitudes in February would be somewhat smaller while the others would remain approximately the same. Though he has not compared these elements with all the observations, a comparison that is also in part entirely unnecessary, he believes he can say in advance that they strike a balance between the observations as far as possible, and that the errors for every observation (apart from those of 3 and 13 January) stay within $5''$ in longitude and latitude, cease to show any regularity and instead change sign quite irregularly. Dr.\ \textsc{Gauss} has good reason to doubt whether with \textit{the very same data} one could obtain \textit{noticeably closer} agreement using other elements. He emphatically does not mean to say or to be understood as saying that other, substantially different elements could not give similar agreement, especially if somewhat different determinations for the Sun were used in this delicate calculation, where even a few seconds make a perceptible difference. Thus \textsc{Piazzi}'s longitude of the Sun in February differs from ours by half a minute, although here, close to quadrature, the influence of the Sun's longitude is smaller than in other positions. Dr.\ \textsc{Gauss} therefore thinks it might be useful to determine the errors in the
\SourceLeaf{208}%
solar tables for these times from very accurate observations and to improve the positions of the Sun accordingly. This \textit{fourth} set of elements is as follows:
\begin{center}
{\small\setlength{\tabcolsep}{2pt}
\begin{tabular}{l r}
Aphelion $\ .\ .\ .\ .\ .\ .\ .\ .\ .\ .$&$326^0\ 27'\ 38''$\\
$\AscendingNode\ .\ .\ .\ .\ .\ .\ .\ .\ .\ .\ .\ .\ .$&$81\quad0\ 44$\\
Inclination $\ .\ .\ .\ .\ .\ .\ .\ .\ .\ .\ .$&$10\ 36\ 57$\\
Logarithm of the semimajor axis $\ .$&$0.4420527$\\
Eccentricity $\ .\ .\ .\ .\ .\ .\ .\ .\ .\ .$&$0.0825017$\\
Epoch 1800 31 Dec. $\ .\ .\ .\ .\ .\ .\ .$&$77^0\ 36'\ 34''$\\[4pt]
\multicolumn{2}{c}{Hence:}\\[3pt]
greatest equation of the centre $\ .\ .\ .\ .$&$9^0\ 27'\ 41''$\\
\multicolumn{2}{l}{mean daily heliocentric tropical motion}\\
$\ .\ .\ .\ .\ .\ .\ .\ .\ .\ .\ .\ .$&$770''914$
\end{tabular}}
\end{center}
From these elements Dr.\ \textsc{Gauss} has calculated in advance the following positions of \textit{Ceres Ferdinandea}. The time is mean time for midnight in \textit{Palermo}.
\begin{center}
{\fontsize{9}{11}\selectfont\setlength{\tabcolsep}{3pt}
\begin{tabular}{r|r|r|r|r|r}
1801&\shortstack{Geocentric\\longitude}&\shortstack{Geocentric\\latitude\\north}&\shortstack{Logarithm\\of the distance\\from $\EarthSymbol$}&\shortstack{Logarithm\\of the distance\\from $\odot$}&\shortstack{Ratio of\\apparent\\brightness}\\\hline
November 25&$5^{\mathrm z}\ 20^0\ 16'$&$9^0\ 25'$&$0.42181$&$0.40468$&$0.6102$\\
December 1&$5\ 22\ 15$&$9\ 48$&$0.40940$&$0.40472$&$0.6459$\\
7&$5\ 24\ 7$&$10\ 12$&$0.39643$&$0.40479$&$0.6855$\\
13&$5\ 25\ 51$&$10\ 37$&$0.38296$&$0.40488$&$0.7290$\\
19&$5\ 27\ 27$&$11\ 4$&$0.36902$&$0.40499$&$0.7770$\\
25&$5\ 28\ 53$&$11\ 32$&$0.35468$&$0.40512$&$0.8295$\\
31&$6\ 0\ 10$&$12\ 1$&$0.34000$&$0.40528$&$0.8869$
\end{tabular}}
\end{center}
If anyone wishes to calculate the position of the planet more precisely from these elements, or for a longer period, we add the following formulas for this purpose:

\begin{center}1) To calculate the equation of the centre $=M$\end{center}
\[
\begin{aligned}
M={}&-34005''494\ \mathrm{Sin.\ Anom.\ med.}\\
&+1750''951\ \mathrm{Sin.}\ 2\ \mathrm{A.\ m.}-124''951\ \mathrm{Sin.}\ 3\ \mathrm{A.\ m}\\
&+10''192\ \mathrm{Sin.}\ 4\ \mathrm{A.\ m.}-0''901\ \mathrm{Sin.}\ 5\ \mathrm{A.\ m.}\\
&+0''083\ \mathrm{Sin.}\ 6\ \mathrm{A.\ m.}
\end{aligned}
\]
\begin{center}2) For the radius vector $=r$\end{center}
\[
\begin{aligned}
r={}&2.776695+0.2373586\ \mathrm{Cosin.\ Anom.\ med.}\\
&-0.0093752\ \mathrm{Cosin.}\ 2\ \mathrm{A.\ m.}\\
&+0.0005789\ \mathrm{Cosin.}\ 3\ \mathrm{A.\ m.}\\
&-0.0000423\ \mathrm{Cosin.}\ 4\ \mathrm{A.\ m.}\\
&+0.0000034\ \mathrm{Cosin.}\ 5\ \mathrm{A.\ m.}\\
&-0.0000003\ \mathrm{Cosin.}\ 6\ \mathrm{A.\ m.}
\end{aligned}
\]
\[
\text{Or:}\quad r=\frac{7.60570}{2.767278\pm0.2283053\ \mathrm{Cosin.\ Anom.\ ver.}}
\]
\begin{center}3) For the heliocentric latitude $=\lambda$.\end{center}
\[
\mathrm{Log.\ Sin.}\ \lambda=9.2653438+\mathrm{Log.\ Sin.\ Arg.\ Latit.}
\]
\begin{center}4) To reduce the heliocentric longitude to the Earth's orbit $=\varepsilon$.\end{center}
\[
\begin{aligned}
\alpha)&\quad\mathrm{Log.\ Tang.}\ \varphi=9.9925025+\mathrm{Log.\ Tang.\ Arg.\ Latit.}\\
\beta)&\quad\varepsilon=\mathrm{Arg.\ Latit.}-\varphi.
\end{aligned}
\]
\begin{center}5) For the reduction of the radius vector, or the shortened distance $=\rho$.\end{center}
\[
\rho=\mathrm{Cosin.\ Latit.\ helioc.}\times\mathrm{Distant.\ vera.}
\]
\begin{center}Or for the logarithm of the shortening itself $=\mathrm{Log.}\ \mu$.\end{center}
\[
\mathrm{Log.}\ \mu=10.000000-\mathrm{Log.\ Cosin.\ Latit.\ helioc.}
\]
\begin{center}6) For the aberration of light in longitude, latitude, right ascension and declination $=a$.\end{center}
\[
\begin{aligned}
\mathrm{Log.}\ a={}&\pm\mathrm{Log.\ Dist.}\ \text{à}\ \EarthSymbol+\mathrm{Log.\ mot.\ horar.\ geoc.}\\
&+\mathrm{Log.}\ 7.751007.
\end{aligned}
\]

\SourceLeaf{209}%
We conclude this report with a remark of Dr.\ \textsc{Gauss} on the inclination of the orbit of \textit{Ceres}, the magnitude of which struck many astronomers. He agrees entirely with some astronomers that no physical reason entitles us to expect all bodies in our solar system with nearly circular orbits also to be slightly inclined to the ecliptic. Yet it seems to him that the paradox would be considerably diminished and that the analogy could still be preserved to some extent if this planet's orbit were referred to its proper reference plane, namely the Sun's equator, as \textsc{La Place} did for the orbits of the satellites of Uranus (\textit{A.\ G.\ E.}\ vol.\ II, p.\ 259). If we compare the orbital planes of all the planets with one another, it then turns out that the orbit of \textit{Ceres} has a greater inclination to the Earth's orbit than to the orbit of any other planet. So when we compare the orbit of \textit{Ceres} with the Earth's orbit, we are comparing the two very extremes of our solar system. If, however, we imagine a plane lying approximately midway among the orbital planes of all \textit{eight} planets, the inclinations to that plane are all small enough. It seems quite remarkable that the Sun's equator itself lies approximately in such a position, with only the qualification that the orbits of \textit{Ceres}, Venus and Mercury are the least inclined to it, by about 3 to 4 degrees, whereas the Earth's orbit is the most inclined.
\begin{flushright}[\textsc{von Zach}.]\end{flushright}


